\documentclass[11pt]{article}

\usepackage[a4paper,margin=1in]{geometry}
\usepackage{amsmath,amssymb}
\usepackage{hyperref}

\hypersetup{
    colorlinks=true,
    linkcolor=blue,
    citecolor=blue,
    urlcolor=blue
}

\title{Directional Derivative}
\author{}
\date{}

\begin{document}

\maketitle

\section*{In Euclidean Space}

Let
\[
f:\mathbb R^n\to\mathbb R
\]
be smooth. The directional derivative of $f$ at $x$ in the direction $v$ is
\[
Df(x)[v]
=
\left.\frac{d}{dt}f(x+tv)\right|_{t=0}.
\]
It measures the first-order rate of change of $f$ if one starts at $x$ and
moves infinitesimally with velocity $v$.

If $f$ is differentiable, then
\[
Df(x)[v]
=
\sum_{i=1}^n
\frac{\partial f}{\partial x^i}(x)v^i.
\]
Thus, for fixed $x$, the map
\[
v\mapsto Df(x)[v]
\]
is linear in $v$.

\section*{On a Smooth Manifold}

Let $M$ be a smooth manifold, let $p\in M$, and let
\[
X_p\in T_pM
\]
be a tangent vector. A tangent vector can be represented by a smooth curve
\[
\gamma:(-\epsilon,\epsilon)\to M
\]
such that
\[
\gamma(0)=p,
\qquad
\dot\gamma(0)=X_p.
\]
The directional derivative of $f:M\to\mathbb R$ at $p$ in the direction $X_p$
is
\[
X_p[f]
=
\left.\frac{d}{dt}f(\gamma(t))\right|_{t=0}.
\]
This number depends only on the tangent vector $X_p$, not on the particular
curve chosen to represent it.

\section*{Relation to the Differential}

The differential of $f$ at $p$ is the covector
\[
df_p:T_pM\to\mathbb R
\]
defined by
\[
df_p(X_p)=X_p[f].
\]
So the differential packages all directional derivatives at $p$ into one
linear map.

\section*{Is It the Covariant Derivative?}

For a scalar function, the directional derivative $X[f]$ is already intrinsic.
It does not require a connection.

A connection is needed when differentiating vector fields or tensor fields
along directions, because values of vector fields at different points live in
different tangent spaces. But a function has values in $\mathbb R$, so along a
curve $\gamma(t)$ we can simply differentiate the real-valued function
\[
t\mapsto f(\gamma(t)).
\]
Therefore:
\[
\text{directional derivative of a function: no connection needed,}
\]
while
\[
\text{covariant derivative of vector/tensor fields: connection needed.}
\]

\end{document}
